% =====================================================================
%  kitap-stil.tex — Sayfa düzeni, yazı tipleri, başlıklar, kutular
%  Lean 4 ile Modern Mantık
% =====================================================================

% ---------- Dil ve yazı tipleri (XeLaTeX) ----------
\usepackage{fontspec}
\usepackage{polyglossia}
\setdefaultlanguage{turkish}
\setotherlanguage{english}

\usepackage{amsmath,amsthm}
\usepackage{unicode-math}
\setmainfont{TeX Gyre Pagella}
\setsansfont{TeX Gyre Heros}[Scale=MatchLowercase]
\setmonofont{DejaVu Sans Mono}[Scale=0.84]
\setmathfont{Latin Modern Math}[Scale=MatchLowercase]
\newfontfamily\leanfont{DejaVu Sans Mono}[Scale=0.84]
\newfontfamily\baslikfont{TeX Gyre Heros}[Ligatures=TeX]
\newfontfamily\sembolfont{DejaVu Sans}   % ● ○ ⚠ ✓ ◇ □ gibi semboller için
\newfontfamily\sembolkodfont{DejaVu Sans}[Scale=0.84]  % kod içinde ⋀ ⋁ ⟹ ⊨

% ---------- Sayfa ----------
\usepackage[a4paper, inner=2.6cm, outer=2.4cm, top=2.6cm, bottom=2.8cm,
            headheight=14pt]{geometry}
\usepackage{microtype}
\usepackage[table]{xcolor}
\usepackage{graphicx}
\usepackage[most]{tcolorbox}
\usepackage{tikz}
\usetikzlibrary{arrows.meta,positioning,calc,shapes.geometric,
                fit,backgrounds,decorations.pathreplacing,matrix}
\usepackage{booktabs,array,tabularx,multirow}
\usepackage{enumitem}
\setlist{itemsep=2pt, topsep=4pt}
\usepackage{ebproof}           % doğal çıkarım ağaçları
\usepackage{csquotes}
\usepackage{multicol}
\usepackage[numbers,sort&compress]{natbib}
\usepackage{xurl}

% ---------- Ana renkler ----------
\definecolor{anarenk}{HTML}{3B5BDB}     % lacivert
\definecolor{ikincil}{HTML}{0CA678}     % yeşil
\definecolor{vurgu}{HTML}{D6336C}       % pembe-kırmızı
\definecolor{koyu}{HTML}{1F2A44}
\definecolor{tabloust}{HTML}{3B5BDB}
\definecolor{tablosatir}{HTML}{EDF2FF}

\definecolor{tanimbg}{HTML}{EDF2FF}  \definecolor{tanimframe}{HTML}{3B5BDB}
\definecolor{teorembg}{HTML}{FFF0F6} \definecolor{teoremframe}{HTML}{C2255C}
\definecolor{ornekbg}{HTML}{F3F0FF}  \definecolor{ornekframe}{HTML}{7950F2}
\definecolor{fikirbg}{HTML}{E6FCF5}  \definecolor{fikirframe}{HTML}{0CA678}
\definecolor{dikkatbg}{HTML}{FFF4E6} \definecolor{dikkatframe}{HTML}{E8590C}
\definecolor{arastirmabg}{HTML}{F1F3F5}\definecolor{arastirmaframe}{HTML}{495057}
\definecolor{tarihbg}{HTML}{FFF9DB}  \definecolor{tarihframe}{HTML}{E67700}
\definecolor{alistirmabg}{HTML}{F8F9FA}\definecolor{alistirmaframe}{HTML}{1971C2}

% ---------- Başlıklar ----------
\usepackage{titlesec}
\titleformat{\chapter}[display]
  {\baslikfont\color{koyu}}
  {\begin{tikzpicture}[remember picture, overlay]
     \fill[anarenk] (current page.north west) rectangle
           ([yshift=-3.2cm]current page.north east);
     \node[anchor=south east, text=white, font=\baslikfont\bfseries\fontsize{60}{60}\selectfont]
          at ([xshift=-2.4cm, yshift=-3.0cm]current page.north east) {\thechapter};
     \node[anchor=south west, text=white!85!anarenk, font=\baslikfont\large]
          at ([xshift=2.6cm, yshift=-3.0cm]current page.north west) {\MakeUppercase{\chaptertitlename}};
   \end{tikzpicture}}
  {0pt}
  {\vspace*{1.2cm}\Huge\bfseries\raggedright}
  [\vspace{4pt}{\color{anarenk}\titlerule[1.5pt]}]
\titlespacing*{\chapter}{0pt}{0pt}{24pt}

\titleformat{\section}
  {\baslikfont\Large\bfseries\color{koyu}}
  {\textcolor{anarenk}{\thesection}}{0.8em}{}
\titleformat{\subsection}
  {\baslikfont\large\bfseries\color{koyu}}
  {\textcolor{anarenk}{\thesubsection}}{0.8em}{}
\titleformat{\subsubsection}
  {\baslikfont\normalsize\bfseries\color{koyu}}
  {\textcolor{anarenk}{\thesubsubsection}}{0.8em}{}
\titleformat{\paragraph}[runin]
  {\baslikfont\normalsize\bfseries\color{koyu}}{}{0pt}{}[.]
\setcounter{secnumdepth}{2}
\setcounter{tocdepth}{1}

\titleformat{\part}[display]
  {\centering\baslikfont\color{white}}
  {\begin{tikzpicture}[remember picture, overlay]
     \fill[koyu] (current page.south west) rectangle (current page.north east);
     \fill[anarenk] ([yshift=-9cm]current page.north west) rectangle
                    ([yshift=-9.4cm]current page.north east);
   \end{tikzpicture}\Large KISIM \thepart}
  {12pt}{\Huge\bfseries}

% ---------- Üst bilgi ----------
\usepackage{fancyhdr}
\pagestyle{fancy}
\fancyhf{}
\renewcommand{\chaptermark}[1]{\markboth{#1}{}}
\renewcommand{\sectionmark}[1]{\markright{\thesection\quad #1}}
\fancyhead[LE]{\baslikfont\small\color{black!60}\ifnum\value{chapter}>0 \thechapter\quad\fi\leftmark}
\fancyhead[RO]{\baslikfont\small\color{black!60}\rightmark}
\fancyfoot[LE,RO]{\baslikfont\small\bfseries\color{anarenk}\thepage}
\renewcommand{\headrulewidth}{0.4pt}
\renewcommand{\headrule}{{\color{black!20}\hrule width\headwidth height\headrulewidth}}
\fancypagestyle{plain}{\fancyhf{}\fancyfoot[C]{\baslikfont\small\color{anarenk}\thepage}
  \renewcommand{\headrulewidth}{0pt}}

% ---------- Numaralı kutular (bölüm bazında) ----------
\newcounter{kutusayac}[chapter]
\renewcommand{\thekutusayac}{\thechapter.\arabic{kutusayac}}

\tcbset{
  kitapkutu/.style={
    enhanced, breakable,
    boxrule=0.6pt, arc=3pt, leftrule=4pt,
    fonttitle=\baslikfont\bfseries,
    titlerule=0pt,
    before skip=12pt, after skip=12pt,
    left=8pt, right=8pt,
  }
}

\newtcbtheorem[number within=chapter]{tanim}{Tanım}{kitapkutu,
  colback=tanimbg, colframe=tanimframe, colbacktitle=tanimbg,
  coltitle=tanimframe!80!black, separator sign={.},
  description font=\mdseries}{tnm}
\newtcbtheorem[use counter from=tanim]{teorem}{Teorem}{kitapkutu,
  colback=teorembg, colframe=teoremframe, colbacktitle=teorembg,
  coltitle=teoremframe!80!black, separator sign={.},
  description font=\mdseries}{trm}
\newtcbtheorem[use counter from=tanim]{ornek}{Örnek}{kitapkutu,
  colback=ornekbg, colframe=ornekframe, colbacktitle=ornekbg,
  coltitle=ornekframe!80!black, separator sign={.},
  description font=\mdseries}{orn}

% Numarasız kutular
\newtcolorbox{onemlifikir}[1][Önemli Fikir]{kitapkutu,
  colback=fikirbg, colframe=fikirframe, colbacktitle=fikirbg,
  coltitle=fikirframe!70!black, title={#1}}
\newtcolorbox{gunlukornek}[1][Günlük Hayattan]{kitapkutu,
  colback=ornekbg, colframe=ornekframe, colbacktitle=ornekbg,
  coltitle=ornekframe!70!black, title={#1}}
\newtcolorbox{dikkat}[1][Dikkat: Sık Yapılan Hata]{kitapkutu,
  colback=dikkatbg, colframe=dikkatframe, colbacktitle=dikkatbg,
  coltitle=dikkatframe!80!black, title={{\sembolfont ⚠}\enspace #1}}
\newtcolorbox{tarihce}[1][Tarihçe]{kitapkutu,
  colback=tarihbg, colframe=tarihframe, colbacktitle=tarihbg,
  coltitle=tarihframe!80!black, title={#1}, fontupper=\small}
\newtcolorbox{arastirmanotu}[1][Araştırma Notu]{kitapkutu,
  colback=arastirmabg, colframe=arastirmaframe, colbacktitle=arastirmaframe,
  coltitle=white, title={#1}, fontupper=\small, titlerule=0.5pt,
  attach boxed title to top left={xshift=10pt,yshift=-2.2mm},
  boxed title style={colback=arastirmaframe, arc=2pt, boxrule=0pt}}

% Bölüm başı: hedefler kutusu
\newtcolorbox{bolumhedefleri}{enhanced, breakable,
  colback=white, colframe=anarenk, boxrule=0.8pt, arc=4pt,
  title={\baslikfont\bfseries Bu bölümde öğrenecekleriniz},
  colbacktitle=anarenk, coltitle=white,
  before skip=6pt, after skip=16pt, left=8pt}

% Bölüm sonu özeti
\newtcolorbox{bolumozeti}{enhanced, breakable,
  colback=fikirbg, colframe=ikincil, boxrule=0.8pt, arc=4pt,
  title={\baslikfont\bfseries Bölüm Özeti},
  colbacktitle=ikincil, coltitle=white,
  before skip=12pt, after skip=12pt, left=8pt}

% ---------- Alıştırmalar ----------
\newcounter{alistirma}[chapter]
\renewcommand{\thealistirma}{\thechapter.\arabic{alistirma}}
\newcommand{\zorluk}[1]{%
  {\sembolfont\footnotesize\ifcase#1\or\textcolor{ikincil}{●○○}\or\textcolor{dikkatframe}{●●○}%
  \or\textcolor{vurgu}{●●●}\fi}}
\newenvironment{alistirmalar}{%
  \begin{tcolorbox}[enhanced, breakable, colback=alistirmabg,
     colframe=alistirmaframe, boxrule=0.6pt, arc=3pt, leftrule=4pt,
     title={\baslikfont\bfseries Alıştırmalar}, colbacktitle=alistirmabg,
     coltitle=alistirmaframe!80!black, titlerule=0pt]
  \small
  Zorluk: \zorluk{1}~temel \quad \zorluk{2}~orta \quad \zorluk{3}~ileri.
  Lean alıştırmalarının iskeletleri eşlik eden projedeki
  \lc{Alistirmalar/} klasöründedir; \leantac{sorry} yerine kendi ispatınızı yazın.
  \begin{description}[leftmargin=0pt, itemsep=6pt, font=\normalfont]}
  {\end{description}\end{tcolorbox}}
\newcommand{\alistirma}[1]{\refstepcounter{alistirma}%
  \item[\baslikfont\bfseries\textcolor{alistirmaframe}{\thealistirma}\enspace\zorluk{#1}]}

% ---------- Doğal çıkarım kuralları ----------
\newcommand{\kural}[1]{\textcolor{vurgu}{\textsf{\small #1}}}
\ebproofset{right label template=\kural{\inserttext}, separation=1.4em}

% ---------- Tablolar ----------
\newcommand{\tablobaslik}[1]{\textcolor{white}{\baslikfont\bfseries\small #1}}

% ---------- Bağlantılar ----------
\usepackage[hidelinks]{hyperref}
\hypersetup{colorlinks=true, linkcolor=anarenk, citecolor=ikincil,
            urlcolor=vurgu, pdfauthor={}, pdftitle={Lean 4 ile Modern Mantık}}
\usepackage[english,nameinlink]{cleveref}  % polyglossia-turkish ile uyum için dil açıkça verilir
\crefname{chapter}{Bölüm}{Bölüm}
\crefname{section}{Kesim}{Kesim}
\crefname{subsection}{Kesim}{Kesim}
\crefname{figure}{Şekil}{Şekil}
\crefname{table}{Tablo}{Tablo}
\crefname{tcb@cnt@tanim}{Tanım}{Tanım}
\crefname{tcb@cnt@teorem}{Teorem}{Teorem}
\crefname{tcb@cnt@ornek}{Örnek}{Örnek}
\crefname{alistirma}{Alıştırma}{Alıştırma}
\crefname{part}{Kısım}{Kısım}

% Şekil/tablo başlıkları
\usepackage[font=small, labelfont={bf,color=anarenk}, labelsep=period]{caption}

% Kısaltmalar
\newcommand{\Lean}{Lean~4}
\newcommand{\Mathlib}{\textsf{Mathlib}}
\newcommand{\terim}[2]{\emph{#1} (İng.\ \foreignlanguage{english}{\emph{#2}})}

% ---------- Taslak bölüm planı ----------
\newtcolorbox{taslakplan}{enhanced, breakable,
  colback=white, colframe=black!35, boxrule=0.8pt, arc=4pt,
  borderline={0.4pt}{2pt}{black!20, dashed},
  title={\baslikfont\bfseries Yazım planı \hfill\normalfont\small taslak},
  colbacktitle=black!8, coltitle=black!70,
  before skip=10pt, after skip=10pt, left=8pt}

% Numaralı adım listeleri:  \begin{enumerate}[adimlar] ... \end{enumerate}
\DeclareRobustCommand{\adimno}[1]{\tikz[baseline=(n.base)]\node[circle,
  fill=anarenk, text=white, inner sep=1.4pt, minimum size=1.15em,
  font=\baslikfont\scriptsize\bfseries](n){#1};}
\SetEnumitemKey{adimlar}{label=\adimno{\arabic*}, leftmargin=2em, itemsep=5pt}

% Boş kalan sol sayfalarda üst/alt bilgi gösterme
\makeatletter
\renewcommand{\cleardoublepage}{\clearpage\if@twoside\ifodd\c@page\else
  \hbox{}\thispagestyle{empty}\newpage\if@twocolumn\hbox{}\newpage\fi\fi\fi}
\makeatother
\colorlet{gri}{black!30}

% Nesne dilinin bağlaçları (Bölüm 10–13): metin içinde
\newcommand{\ove}{\mathbin{\textstyle\bigwedge}}
\newcommand{\oveya}{\mathbin{\textstyle\bigvee}}
\newcommand{\oise}{\Longrightarrow}

% ---- Cilt II'ye göndermeler ----
% Kullanım: \ciltiki{4}'te  →  "Cilt II, Bölüm 4'te"
\newcommand{\ciltiki}[1]{Cilt~II, Bölüm~#1}

% ---- Proje kartları (Bölüm 17) ----
\newtcolorbox{projekarti}[2]{kitapkutu,
  colback=white, colframe=ornekframe, colbacktitle=ornekframe,
  coltitle=white, title={#1\enspace\textbar\enspace #2}}
\newcommand{\projebilgi}[3]{%
  \noindent{\baslikfont\small
    \textcolor{ornekframe}{\bfseries Ön koşul:}~#1\hfill
    \textcolor{ornekframe}{\bfseries Zorluk:}~#2\hfill
    \textcolor{ornekframe}{\bfseries Tahmini emek:}~#3}\par\medskip}
